\documentclass[a4paper,12pt]{article}
\usepackage[utf8]{inputenc}
\usepackage[ngerman]{babel}
\usepackage{amsmath,amssymb}
\usepackage{geometry}
\geometry{margin=2.5cm}

\title{Aufgabe 7(d) -- Gauß- und Gauß-Jordan-Elimination}
\author{}
\date{}

\begin{document}
\maketitle

\section*{Matrixgleichung}

\[
\begin{pmatrix} 1 & 0 & 0 & 0 & 3 \\ 0 & 1 & 0 & 2 & 0 \\ 0 & 1 & 0 & 1 & 0 \\ 2 & 3 & 0 & 4 & 6 \end{pmatrix}
\begin{pmatrix} x \\ y \\ z \\ t \\ w \end{pmatrix}
=
\begin{pmatrix} 2 \\ 2 \\ 1 \\ 8 \end{pmatrix}
\]

Das entspricht dem Gleichungssystem:
\[
\begin{cases}
x + 3w = 2 \\
y + 2t = 2 \\
y + t = 1 \\
2x + 3y + 4t + 6w = 8
\end{cases}
\]

\section*{Erweiterte Koeffizientenmatrix}

\[
\left(\begin{array}{ccccc|c}
1 & 0 & 0 & 0 & 3 & 2 \\
0 & 1 & 0 & 2 & 0 & 2 \\
0 & 1 & 0 & 1 & 0 & 1 \\
2 & 3 & 0 & 4 & 6 & 8
\end{array}\right)
\]

\section*{Gauß-Elimination (Stufenform)}

\textbf{Schritt 1:} $Z_4 \leftarrow Z_4 - 2Z_1$

\[
\left(\begin{array}{ccccc|c}
1 & 0 & 0 & 0 & 3 & 2 \\
0 & 1 & 0 & 2 & 0 & 2 \\
0 & 1 & 0 & 1 & 0 & 1 \\
0 & 3 & 0 & 4 & 0 & 4
\end{array}\right)
\]

\textbf{Schritt 2:} $Z_3 \leftarrow Z_3 - Z_2$ und $Z_4 \leftarrow Z_4 - 3Z_2$

\[
\left(\begin{array}{ccccc|c}
1 & 0 & 0 & 0 & 3 & 2 \\
0 & 1 & 0 & 2 & 0 & 2 \\
0 & 0 & 0 & -1 & 0 & -1 \\
0 & 0 & 0 & -2 & 0 & -2
\end{array}\right)
\]

\textbf{Schritt 3:} $Z_4 \leftarrow Z_4 - 2Z_3$

\[
\left(\begin{array}{ccccc|c}
1 & 0 & 0 & 0 & 3 & 2 \\
0 & 1 & 0 & 2 & 0 & 2 \\
0 & 0 & 0 & -1 & 0 & -1 \\
0 & 0 & 0 & 0 & 0 & 0
\end{array}\right)
\]

\textbf{Ergebnis:} 3 Pivot-Elemente in den Spalten 1, 2 und 4. Keine Widerspruchszeile.

Freie Variablen: $z$ und $w$ (Spalten 3 und 5 haben kein Pivot).

Das System ist \textbf{konsistent} mit \textbf{unendlich vielen Lösungen}.

\textbf{Bemerkung:} Die Variable $z$ kommt in keiner Gleichung vor (Spalte 3 besteht nur aus Nullen), sie ist also komplett frei.

\subsection*{Rückwärtseinsetzen}

Aus Zeile 3: $-t = -1 \Rightarrow t = 1$

Aus Zeile 2: $y + 2t = 2 \Rightarrow y = 2 - 2 \cdot 1 = 0$

Aus Zeile 1: $x + 3w = 2 \Rightarrow x = 2 - 3w$

\section*{Gauß-Jordan-Elimination (reduzierte Stufenform)}

Wir starten von der Stufenform und normieren/eliminieren nach oben.

\textbf{Schritt 4:} $Z_3 \leftarrow -Z_3$ (Pivot normieren)

\[
\left(\begin{array}{ccccc|c}
1 & 0 & 0 & 0 & 3 & 2 \\
0 & 1 & 0 & 2 & 0 & 2 \\
0 & 0 & 0 & 1 & 0 & 1 \\
0 & 0 & 0 & 0 & 0 & 0
\end{array}\right)
\]

\textbf{Schritt 5:} $Z_2 \leftarrow Z_2 - 2Z_3$ (eliminiere $t$ aus Zeile 2)

\[
\left(\begin{array}{ccccc|c}
1 & 0 & 0 & 0 & 3 & 2 \\
0 & 1 & 0 & 0 & 0 & 0 \\
0 & 0 & 0 & 1 & 0 & 1 \\
0 & 0 & 0 & 0 & 0 & 0
\end{array}\right)
\]

Zeile 1 hat bereits 0 in den Pivot-Spalten 2 und 4, also ist dies die \textbf{reduzierte Zeilenstufenform} (RREF).

\section*{Allgemeine Lösung}

\textbf{Pivot-Variablen:} $x, y, t$ \quad \textbf{Freie Variablen:} $z, w$

Setze $z = s$, $w = q$ mit $s, q \in \mathbb{R}$.

\subsection*{Closed Form (als $n$-Tupel)}

\[
\boxed{(x,y,z,t,w) = (2-3q,\; 0,\; s,\; 1,\; q), \quad s, q \in \mathbb{R}}
\]

\subsection*{Expanded Form (Vektordarstellung)}

\[
\begin{pmatrix} x \\ y \\ z \\ t \\ w \end{pmatrix}
=
\underbrace{\begin{pmatrix} 2 \\ 0 \\ 0 \\ 1 \\ 0 \end{pmatrix}}_{v_0}
+ s \underbrace{\begin{pmatrix} 0 \\ 0 \\ 1 \\ 0 \\ 0 \end{pmatrix}}_{v_1}
+ q \underbrace{\begin{pmatrix} -3 \\ 0 \\ 0 \\ 0 \\ 1 \end{pmatrix}}_{v_2}
\]

\[
\boxed{\vec{x} = v_0 + s \cdot v_1 + q \cdot v_2, \quad s, q \in \mathbb{R}}
\]

\section*{Probe}

Setze z.\,B.\ $s = 0, q = 0$, also $(x,y,z,t,w) = (2,0,0,1,0)$:

\begin{align*}
2 + 3 \cdot 0 &= 2 \checkmark \\
0 + 2 \cdot 1 &= 2 \checkmark \\
0 + 1 &= 1 \checkmark \\
2 \cdot 2 + 3 \cdot 0 + 4 \cdot 1 + 6 \cdot 0 &= 4 + 4 = 8 \checkmark
\end{align*}

\end{document}
